% 第 6 题：理想变压器电路（circuitikz）
% 交流电源—R2（并联电压表）—原线圈；副线圈经 P1 抽头接 R1
\begin{tikzpicture}[font=\small, >=Stealth]
  % 交流电源
  \draw (0,0) to[sV] (0,3);
  % 上边：滑动变阻器 R2
  \draw (0,3) -- (1.6,3);
  \draw (1.6,3) to[vR, l=$R_2$] (4.6,3);
  \draw (4.6,3) -- (5,3);
  \fill (1.6,3) circle (1.5pt);
  \fill (4.6,3) circle (1.5pt);
  % 电压表并联在 R2 两端
  \draw (1.6,3) -- (1.6,1.8) to[voltmeter] (4.6,1.8) -- (4.6,3);
  % 下边：电流表
  \draw (0,0) -- (2.9,0);
  \draw (2.9,0) to[ammeter] (3.9,0);
  \draw (3.9,0) -- (5,0);
  % 原线圈
  \draw (5,3) -- (5,2.6);
  \draw (5,2.6) to[L] (5,0.4);
  \draw (5,0.4) -- (5,0);
  % 铁芯
  \draw[thick] (5.18,0.15) -- (5.18,2.85);
  \draw[thick] (5.34,0.15) -- (5.34,2.85);
  % 副线圈
  \draw (5.75,2.6) to[L] (5.75,0.4);
  % P1 抽头与 R1
  \draw (5.75,1.5) -- (7.1,1.5);
  \draw (7.1,1.5) to[R, l=$R_1$] (7.1,0.4);
  \draw (7.1,0.4) -- (5.75,0.4);
  \fill (5.75,1.5) circle (1.5pt);
  \draw[->] (6.55,1.95) -- (5.8,1.55);
  \node at (6.65,2.1) {$\mathrm{P}_1$};
  \node at (2.6,3.4) {$\mathrm{P}_2$};
\end{tikzpicture}
